\section{AlphaProof: RL se encuentra con las matemáticas formales}

\subsection{Lean como el entorno de RL ideal}
Lean ofrece una tupla completa de state/action/reward/observation. Hubert mostró un episode típico: a partir del `goal` se abstrae un lemma graph → el agent elige `apply`/`intro`/`rw` → Lean intenta simplificar el goal → éxito o fail. «Al final tratamos la salida de Lean como el simulator del entorno de RL», dijo.
\begin{importantbox}{Anatomía del entorno}
1) Action space: tactics en capas (primitive tactic, combination tactic, search macro); 2) Observation: árbol de goal + local context + type info; 3) Reward: binario; 4) Reset: en cada new theorem se reinicializa el árbol de MCTS.
\end{importantbox}

\subsection{Arquitectura del sistema y pipeline de entrenamiento}
AlphaProof está compuesto por varias redes: el Formalizer traduce proposiciones en lenguaje natural/LaTeX en aserciones de Lean; el Proof Generator/Actor-critic se encarga de la search; el Verifier realiza una segunda verificación; el Runner registra el trace.
\begin{table}[H]
\centering
\begin{tabular}{p{4.5cm}p{4.5cm}p{4cm}}
\toprule
Componente & Responsabilidad & Notas \\
\midrule
Formalizer & Estructura las proposiciones, genera declarations de Lean & Basado en seq-to-seq + retrieval \\
Proof Generator & RL + search genera tactic sequence & Combinación de Monte Carlo y beam \\
Verifier & Lean verifica el proof term & Si falla, dispara un rollback \\
Runner & Registra proof trace + metrics & Se usa para el governance gate \\
\bottomrule
\end{tabular}
\caption{Componentes de la arquitectura de AlphaProof}
\end{table}
\begin{knowledgebox}{Puntos de gobernanza del pipeline}
Cada component escribe en el governance log: la comparación de semantic gap del Formalizer, la entropy/confidence del Proof Generator, la failure rate del Verifier, el chunk gating hit del Runner. Hubert dijo: «Solo si convertimos el trace en un artifact estructurado, el RL puede dejar de ser una caja negra».
\end{knowledgebox}

Durante el entrenamiento, AlphaProof alterna el entrenamiento entre varios dataset: Mathlib core definitions, IMO contest statements, AlphaGeometry proofs. Hubert enfatizó que cada dataset corresponde a un replay buffer independiente (el «Lemma buffer»), y que en el training schedule se establecen 3 etapas, warm start → search-intensive → human eval, lo que hace que la policy sea cada vez más robusta.

\subsection{Experimentos, IMO 2024 e iteración}
AlphaProof, junto con AlphaGeometry, obtuvo la medalla de plata en el IMO 2024; los principales avances fueron: 1) construir un large-scale formalization dataset; 2) usar el replay buffer para registrar el failure trace durante la proof search; 3) usar el resultado del human review del governance board como reward shaping. Hubert mostró en particular un ejemplo: cuando el agent tiene un fail en el teorema `prime`, el trace se sube a `AlphaProof Repos`, un math expert genera un corrected lemma, y este se feed de nuevo a la policy.
\begin{warningbox}{Advertencias para el despliegue a nivel IMO}
1) hay que probar de forma segmentada parámetros como search depth, tactic mix, chunk gating; 2) cuando el tiempo de permanencia del proof supera los 500 ms hay que forzar el log; 3) no basta con mirar el final success rate, también hay que hacer tally de los drift gate triggers.
\end{warningbox}

\subsection{Resumen de la sección}
AlphaProof aprovecha el deterministic reward, structured observation y rich action space que ofrece Lean para convertir la math proof search en un RL pipeline controlable; el log de gobernanza y el human review hacen que incluso el fracaso sea útil.
\newpage

